% Lecture notes: Week 13, Lecture 3 -- Thin-Walled Pressure Vessels
% To hide this lecture from the website, add a line containing only:  % draft
% To set the date shown on the website, add a line like:  % posted: 2026-09-02
\documentclass[11pt]{article}
\usepackage{mech2030notes}
\begin{document}

% ##################################################################
%  WEEK 13, LECTURE 3
% ##################################################################
\lecture{13}{3}{Thin-Walled Pressure Vessels}

\noindent
(The summary below is what was presented in lecture. The full derivations are given
afterwards, in \hyperref[w13l3:deriv]{Detailed derivations}.)

% ==================================================================
\section*{Summary}
% ==================================================================

\subsection*{Cylindrical vessel}

\noindent
\begin{minipage}[c]{0.62\textwidth}
\centering
\begin{tikzpicture}[scale=0.95]
  % 3D triad
  \begin{scope}[shift={(-1.9,0.2)}]
    \draw[axis] (0,0) -- (0,1.0) node[above] {$z$};
    \draw[axis] (0,0) -- (1.0,0) node[right] {$y$};
    \draw[axis] (0,0) -- (-0.6,-0.6) node[below left=-2pt] {$x$};
  \end{scope}
  % cylinder body
  \fill[gray!10] (0,-1.1) -- (4.6,-1.1) arc[start angle=-90, end angle=90,
       x radius=0.35, y radius=1.1] -- (0,1.1) arc[start angle=90, end angle=270,
       x radius=0.35, y radius=1.1] -- cycle;
  \draw[thick] (0,1.1) -- (4.6,1.1);
  \draw[thick] (0,-1.1) -- (4.6,-1.1);
  \draw[thick] (0,1.1) arc[start angle=90, end angle=270, x radius=0.35, y radius=1.1];
  \draw[thick, dashed, gray] (0,-1.1) arc[start angle=-90, end angle=90,
       x radius=0.35, y radius=1.1];
  \draw[thick] (4.6,0) ellipse[x radius=0.35, y radius=1.1];
  % radius
  \draw[thinarrow] (4.6,0) -- (4.6,1.08);
  \fill (4.6,0) circle (0.03);
  \node[left] at (4.6,0.55) {$r$};
  \node at (3.7,-0.75) {$p$};
  % stress element
  \draw[thick, fill=white] (1.95,-0.2) rectangle (2.35,0.2);
  \draw[force] (2.15,0.2) -- (2.15,0.75) node[above] {$\sigma_h$};
  \draw[force] (2.15,-0.2) -- (2.15,-0.75) node[below] {$\sigma_h$};
  \draw[force] (2.35,0) -- (2.95,0) node[right] {$\sigma_\ell$};
  \draw[force] (1.95,0) -- (1.35,0) node[left] {$\sigma_\ell$};
\end{tikzpicture}
\end{minipage}%
\hfill
\begin{minipage}[c]{0.34\textwidth}
\begin{itemize}[leftmargin=1.2em, itemsep=1pt]
  \item radius $r$
  \item wall thickness $t$
  \item \emph{gauge} pressure $p$
\end{itemize}
\end{minipage}

\medskip
\begin{itemize}[label=$\rightarrow$, itemsep=2pt]
  \item $\sigma_h$ is the circumferential ``\emph{hoop stress}.''
  \item $\sigma_\ell$ is the ``\emph{longitudinal stress}.''
  \item The stress in the through-thickness direction is \emph{negligible} if the wall is
        \emph{thin} ($r/t \ge 10$).
\end{itemize}

\noindent
$\rightarrow$ It can be shown that (see the
\hyperref[w13l3:deriv]{detailed derivations} below)
\begin{center}
\fbox{\fbox{$\displaystyle
  \sigma_h = \frac{pr}{t}, \qquad \sigma_\ell = \frac{pr}{2t}
$}}
\qquad
\begin{tabular}{@{}l@{}}
  thin-walled cylindrical \\ pressure vessel \\ ($r/t \ge 10$)
\end{tabular}
\end{center}

\begin{itemize}[label=$\rightarrow$, itemsep=2pt]
  \item Importantly, the \emph{hoop} stress is \emph{twice} the \emph{longitudinal}
        stress.
  \item Also, since $\sigma_t \approx 0$, the stress state is 2D, with
        \emph{principal stresses}
\end{itemize}
\begin{center}
\fbox{\fbox{$\displaystyle
  \sigma_1 = \sigma_h = \frac{pr}{t}, \qquad \sigma_2 = \sigma_\ell = \frac{pr}{2t}
$}}
\end{center}

\subsection*{Spherical vessel}

\noindent
\begin{minipage}[c]{0.36\textwidth}
\centering
\begin{tikzpicture}[scale=1.15]
  \fill[gray!10] (0,0) circle (1.4);
  \draw[thick] (0,0) circle (1.4);
  % equator and meridian
  \draw[semithick] (-1.4,0) arc[start angle=180, end angle=360, x radius=1.4, y radius=0.35];
  \draw[semithick, dashed, gray] (1.4,0) arc[start angle=0, end angle=180, x radius=1.4, y radius=0.35];
  \draw[semithick] (0,1.4) arc[start angle=90, end angle=270, x radius=0.35, y radius=1.4];
  \draw[semithick, dashed, gray] (0,-1.4) arc[start angle=-90, end angle=90, x radius=0.35, y radius=1.4];
  % stress element
  \draw[thick, fill=white] (-0.2,-0.2) rectangle (0.2,0.2);
  \draw[force] (0,0.2) -- (0,0.75) node[above, fill=gray!10, inner sep=1pt] {$\sigma_h$};
  \draw[force] (0,-0.2) -- (0,-0.75) node[below, fill=gray!10, inner sep=1pt] {$\sigma_h$};
  \draw[force] (0.2,0) -- (0.8,0) node[right, fill=gray!10, inner sep=1pt] {$\sigma_\ell$};
  \draw[force] (-0.2,0) -- (-0.8,0) node[left, fill=gray!10, inner sep=1pt] {$\sigma_\ell$};
\end{tikzpicture}
\end{minipage}%
\hfill
\begin{minipage}[c]{0.6\textwidth}
$\rightarrow$ In this case, $\sigma_h = \sigma_\ell$ and $\sigma_t \approx 0$:
\begin{center}
\fbox{\fbox{$\displaystyle
  \sigma_1 = \sigma_h = \frac{pr}{2t}, \qquad \sigma_2 = \sigma_\ell = \frac{pr}{2t}
$}}

\smallskip
thin-walled spherical pressure vessel ($r/t \ge 10$)
\end{center}
\end{minipage}

% ==================================================================
\clearpage
\phantomsection\label{w13l3:deriv}
\section*{Detailed derivations}
% ==================================================================

\noindent
(Supplementary material for the online notes.)

\subsection*{Cylindrical vessel: hoop stress}

Consider a cylinder of radius $r$ and wall thickness $t$, under gauge pressure $p$, with
coordinates $x$, $y$ (along the axis), and $z$. Cut a slice of length $dy$ out of the
cylinder, and then cut that slice in half with the $yz$-plane. The free body is the half
ring \emph{plus} the pressurized fluid it contains:

\begin{center}
\begin{tikzpicture}[scale=0.95]
  % fluid (half disk)
  \fill[blue!6] (0,-1.6) arc[start angle=-90, end angle=90, radius=1.6] -- cycle;
  % wall (half annulus)
  \fill[gray!15] (0,1.6) arc[start angle=90, end angle=-90, radius=1.6]
                 -- (0,-1.95) arc[start angle=-90, end angle=90, radius=1.95] -- cycle;
  \draw[thick] (0,1.6) arc[start angle=90, end angle=-90, radius=1.6];
  \draw[thick] (0,1.95) arc[start angle=90, end angle=-90, radius=1.95];
  \draw[thick] (0,1.95) -- (0,-1.95);
  % pressure on the cut face of the fluid
  \foreach \zz in {-1.3,-0.9,-0.5,-0.1,0.3,0.7,1.1}
    \draw[thinarrow] (-0.65,\zz+0.1) -- (-0.03,\zz+0.1);
  \node[left] at (-0.7,0) {$p$};
  % hoop stress on the cut wall faces
  \draw[force] (0,1.775) -- (-0.9,1.775) node[left] {$\sigma_h$};
  \draw[force] (0,-1.775) -- (-0.9,-1.775) node[left] {$\sigma_h$};
  % dimensions
  \fill (0,0) circle (0.035);
  \draw[thinarrow] (0,0) -- (35:1.58);
  \node at (35:0.95) [above left=-3pt] {$r$};
  \draw[dim] (2.25,1.6) -- (2.25,1.95);
  \node[right] at (2.3,1.775) {$t$};
  \draw[dim] (2.25,-1.6) -- (2.25,-1.95);
  \node[right] at (2.3,-1.775) {$t$};
  % triad
  \begin{scope}[shift={(3.4,-1.5)}]
    \draw[axis] (0,0) -- (0.8,0) node[right] {$x$};
    \draw[axis] (0,0) -- (0,0.8) node[above] {$z$};
    \draw[semithick] (0,0) circle (0.1);
    \draw[semithick] (-0.07,-0.07) -- (0.07,0.07) (-0.07,0.07) -- (0.07,-0.07);
    \node[below left] at (0,-0.05) {$y$};
  \end{scope}
  \node[align=center, font=\small] at (3.9,0.4) {slice length $dy$\\(into page)};
\end{tikzpicture}
\end{center}

\noindent
The pressure acts over the projected rectangle $(2r)(dy)$, and the hoop stress acts over
the two cut wall faces, each of area $t\,dy$:
\begin{align*}
  \textstyle\sum F_x = 0
  \quad\rightarrow\quad
  -p\,(2r)\,dy + 2\,\sigma_h\,(t)\,dy = 0
  \quad\Rightarrow\quad
  \sigma_h = \dfrac{2pr}{2t}
\end{align*}
\begin{center}
\fbox{\fbox{$\displaystyle \sigma_h = \frac{pr}{t}$}}
\qquad (hoop stress)
\end{center}

\subsection*{Cylindrical vessel: longitudinal stress}

Now cut the vessel with the $xz$-plane. The free body is the closed end of the vessel plus
the fluid it contains (side view):

\begin{center}
\begin{tikzpicture}[scale=0.95]
  % fluid
  \fill[blue!6] (0,-1.2) rectangle (3.6,1.2);
  % walls: top, bottom, end cap
  \fill[gray!15] (-0.3,1.2) rectangle (3.6,1.5);
  \fill[gray!15] (-0.3,-1.5) rectangle (3.6,-1.2);
  \fill[gray!15] (-0.3,-1.2) rectangle (0,1.2);
  \draw[thick] (3.6,1.5) -- (-0.3,1.5) -- (-0.3,-1.5) -- (3.6,-1.5);
  \draw[thick] (3.6,1.2) -- (0,1.2) -- (0,-1.2) -- (3.6,-1.2);
  \draw[thick] (3.6,1.5) -- (3.6,-1.5);
  % pressure on the fluid cut face
  \foreach \zz in {-0.95,-0.57,-0.19,0.19,0.57,0.95}
    \draw[thinarrow] (4.25,\zz) -- (3.63,\zz);
  \node[right] at (4.3,0) {$p$};
  % longitudinal stress on cut wall faces
  \draw[force] (3.6,1.35) -- (4.6,1.35) node[right] {$\sigma_\ell$};
  \draw[force] (3.6,-1.35) -- (4.6,-1.35) node[right] {$\sigma_\ell$};
  % dimensions
  \draw[dim] (6.0,1.2) -- (6.0,1.5);   \node[right] at (6.05,1.35) {$t$};
  \draw[dim] (6.0,-1.2) -- (6.0,-1.5); \node[right] at (6.05,-1.35) {$t$};
  \draw[dim] (6.0,-1.2) -- (6.0,1.2);  \node[right] at (6.05,0) {$2r$};
  % triad
  \begin{scope}[shift={(-1.6,-0.5)}]
    \draw[axis] (0,0) -- (0.8,0) node[right] {$y$};
    \draw[axis] (0,0) -- (0,0.8) node[above] {$z$};
    \draw[semithick] (0,0) circle (0.1);
    \fill (0,0) circle (0.03);
    \node[below left] at (0,-0.05) {$x$};
  \end{scope}
\end{tikzpicture}
\end{center}

\noindent
The pressure acts on the circular fluid face, of area $\pi r^2$, and the longitudinal
stress acts on the ring-shaped wall face, of area $\approx 2\pi r t$ (thin wall):
\begin{align*}
  \textstyle\sum F_y = 0
  \quad\rightarrow\quad
  -p\,(\pi r^2) + \sigma_\ell\,(2\pi r t) = 0
\end{align*}
\begin{center}
\fbox{\fbox{$\displaystyle \sigma_\ell = \frac{pr}{2t}$}}
\qquad (longitudinal stress)
\end{center}

\subsection*{Thin-walled assumption}

At the inner wall there is also a stress in the thickness direction, since the fluid
pushes on the wall:

\begin{center}
\begin{tikzpicture}
  \fill[gray!15] (0,0) rectangle (1.1,0.7);
  \draw[thick] (0,0) rectangle (1.1,0.7);
  \foreach \zz in {0.15,0.35,0.55}
    \draw[thinarrow] (-0.8,\zz) -- (-0.03,\zz);
  \node[left] at (-0.85,0.35) {$p$};
  \foreach \zz in {0.15,0.35,0.55}
    \draw[thinarrow] (1.9,\zz) -- (1.13,\zz);
  \node[right] at (1.95,0.35) {$\sigma_t$};
  \node[font=\small, gray] at (-0.4,-0.3) {fluid side};
\end{tikzpicture}
\end{center}

\noindent
$|\sigma_t| = p$ (compressive) at the inner wall, decaying to zero at the outer wall.

\medskip\noindent
\then\ However, for \emph{thin-walled} pressure vessels ($r/t \ge 10$),
\[
  \left.
  \begin{aligned}
    |\sigma_t| &= p \\
    \sigma_h = 2\sigma_\ell &= p\left(\frac{r}{t}\right)
  \end{aligned}
  \ \right\}
  \quad\rightarrow\quad
  \boxed{\sigma_h,\ \sigma_\ell \gg |\sigma_t|}
\]
\then\ and as a result, $\sigma_t$ is usually \emph{ignored}.

\subsection*{Principal stresses}

$\rightarrow$ If $\sigma_t \approx 0$, then the hoop and longitudinal stresses are the
principal stresses (there is no shear on the $h$--$\ell$ element):
\begin{center}
\fbox{\fbox{$\displaystyle
  \sigma_1 = \sigma_h = \frac{pr}{t}, \qquad \sigma_2 = \sigma_\ell = \frac{pr}{2t}
$}}
\end{center}

\subsection*{Spherical vessel}

Cut the sphere through its center. By symmetry every such cut looks the same, so
$\sigma_h = \sigma_\ell$. The free body is the hemisphere plus its fluid:

\begin{center}
\begin{tikzpicture}[scale=0.95]
  % fluid
  \fill[blue!6] (0,-1.4) arc[start angle=270, end angle=90, radius=1.4] -- cycle;
  % wall
  \fill[gray!15] (0,1.4) arc[start angle=90, end angle=270, radius=1.4]
                 -- (0,-1.7) arc[start angle=270, end angle=90, radius=1.7] -- cycle;
  \draw[thick] (0,1.4) arc[start angle=90, end angle=270, radius=1.4];
  \draw[thick] (0,1.7) arc[start angle=90, end angle=270, radius=1.7];
  \draw[thick] (0,1.7) -- (0,-1.7);
  % pressure
  \foreach \zz in {-1.0,-0.6,-0.2,0.2,0.6,1.0}
    \draw[thinarrow] (0.65,\zz) -- (0.03,\zz);
  \node[right] at (0.7,0) {$p$};
  % stresses
  \draw[force] (0,1.55) -- (1.0,1.55) node[right] {$\sigma_h = \sigma_\ell$};
  \draw[force] (0,-1.55) -- (1.0,-1.55) node[right] {$\sigma_h = \sigma_\ell$};
  % radius
  \fill (0,0) circle (0.035);
  \draw[thinarrow] (0,0) -- (215:1.38);
  \node at (215:0.8) [below right=-2pt] {$r$};
\end{tikzpicture}
\end{center}

\begin{align*}
  \textstyle\sum F = 0
  \quad\rightarrow\quad
  -p\,(\pi r^2) + \sigma_h\,(2\pi r t) = 0
\end{align*}
\begin{center}
\fbox{\fbox{$\displaystyle
  \sigma_1 = \sigma_h = \frac{pr}{2t}, \qquad \sigma_2 = \sigma_\ell = \frac{pr}{2t}
$}}
\end{center}

\end{document}
